\documentclass{article}
\usepackage[utf8]{inputenc}
\usepackage[spanish]{babel}
\usepackage{listings}
\usepackage{xcolor}
\usepackage{geometry}
\geometry{a4paper, margin=1in}

\% Configuración para código C++
\definecolor{codegreen}{rgb}{0,0.6,0}
\definecolor{codegray}{rgb}{0.5,0.5,0.5}
\definecolor{codepurple}{rgb}{0.58,0,0.82}
\definecolor{backcolour}{rgb}{0.95,0.95,0.92}

\lstset{
    backgroundcolor=\color{backcolour},   
    commentstyle=\color{codegreen},
    keywordstyle=\color{magenta},
    numberstyle=\tiny\color{codegray},
    stringstyle=\color{codepurple},
    basicstyle=\ttfamily\small,
    breakatwhitespace=false,         
    breaklines=true,                 
    captionpos=b,                    
    keepspaces=true,                 
    numbers=left,                    
    numbersep=5pt,                  
    showspaces=false,                
    showstringspaces=false,
    showtabs=false,                  
    tabsize=2,
    language=C++,
    inputencoding=utf8,
    extendedchars=true,
    literate={á}{{\'a}}1 {é}{{\'e}}1 {í}{{\'i}}1 {ó}{{\'o}}1 {ú}{{\'u}}1 {ñ}{{\~n}}1 {¿}{{\textquestiondown}}1 {¡}{{\textexclamdown}}1
}

\% Macros para las secciones de los apuntes
\newcommand{\quees}[1]{
    \section*{🤔 ¿QUÉ ES?}
    \hrulefill \\
    \#1
}

\newcommand{\dichodeotraforma}[1]{
    \section*{💬 DICHO DE OTRA FORMA}
    \hrulefill \\
    \#1
}

\newcommand{\diagrama}[1]{
    \section*{🎨 DIAGRAMA}
    \hrulefill \\
    \#1
}

\newcommand{\comofunciona}[1]{
    \section*{⚙️ ¿CÓMO FUNCIONA?}
    \hrulefill \\
    \#1
}

\newcommand{\complejidad}[1]{
    \section*{📱 COMPLEJIDAD (Big-O)}
    \hrulefill \\
    \#1
}

\newcommand{\entrevistas}[1]{
    \section*{🎯 EN ENTREVISTAS}
    \hrulefill \\
    \#1
}

\newcommand{\resumen}[1]{
    \section*{📋 RESUMEN RÁPIDO}
    \hrulefill \\
    \#1
}

\begin{document}

\title{Complejidad Espacial (Space Complexity)}
\author{Aldo Zepeda Montoya}
\date{23 de mayo de 2026}
\maketitle

\subsection*{CATEGORÍA: 00\_Fundamentos}

\quees{
Imagina que tu programa es una receta de cocina 🍳. Big-O en tiempo
mide cuánto TARDAS en cocinar. Complejidad espacial mide cuántas
OLLAS Y SARTENES necesitas mientras cocinas.

Si para preparar una ensalada de 10 ingredientes necesitas una tabla
de picar pequeña, pero para 100 ingredientes necesitas una mesa gigante,
eso significa que el espacio crece con el número de ingredientes.

En programación, la complejidad espacial mide cuánta MEMORIA EXTRA
(RAM) utiliza tu algoritmo para resolver un problema a medida que el
input se hace más grande.
}

\dichodeotraforma{
La Space Complexity analiza el uso de recursos de memoria. Es crucial
distinguir entre:

\begin{enumerate}
    \item \textbf{Input Space:} La memoria para guardar los datos de entrada.
    \item \textbf{Auxiliary Space:} La memoria EXTRA o temporal que el algoritmo crea.
\end{enumerate}

En la mayoría de las entrevistas, cuando preguntan "Space Complexity",
se refieren al Auxiliary Space. Si tu algoritmo modifica el input original
sin crear estructuras nuevas, se dice que es "In-place" y su complejidad
espacial es O(1).
}

\diagrama{
Cuando usas recursión, cada llamada ocupa un "piso" en el stack:

\begin{verbatim}
factorial(4)
  |-- factorial(3)
       |-- factorial(2)
            |-- factorial(1)  <-- Tope del stack
\end{verbatim}
Cada nivel de recursión consume memoria adicional en el Call Stack.
}

\comofunciona{
\begin{description}
    \item[O(1) --- Espacio Constante] Solo usas unas cuantas variables fijas sin importar el input. Ejemplo: Hacer un swap de dos números usando una variable temporal.
    \item[O(n) --- Espacio Lineal] Creas una estructura (como un array o vector) proporcional al input. Ejemplo: Copiar una lista de n elementos a una nueva lista.
    \item[O(n^2) --- Espacio Cuadrático] Creas una estructura bidimensional (matriz) de n x n. Ejemplo: Una tabla de adyacencia para un grafo denso.
\end{description}
}

\complejidad{
\begin{tabular}{|l|l|l|}
\hline
Algoritmo \& Espacio \& Razón \\
\hline
Acceso a Array \& O(1) \& No requiere memoria extra \\
Búsqueda Lineal \& O(1) \& Solo usa un par de índices \\
Búsqueda Binaria iter. \& O(1) \& Solo usa izq/der/mid \\
Búsqueda Binaria recur. \& O(log n) \& Log(n) llamadas en el stack \\
Merge Sort \& O(n) \& Necesita un array auxiliar \\
Quick Sort \& O(log n) \& Promedio de llamadas stack \\
Crear matriz n x n \& O(n^2) \& Espacio bidimensional \\
\hline
\end{tabular}
}

\entrevistas{
\begin{itemize}
    \item En Google preguntan SIEMPRE espacio Y tiempo. No esperes a que pregunten.
    \item \textbf{"In-place algorithm":} Significa que el algoritmo tiene O(1) espacio auxiliar porque modifica el input original directamente.
    \item \textbf{Trade-off clásico:} A veces puedes hacer que un programa sea más rápido (tiempo) si estás dispuesto a usar más memoria (espacio). Ejemplo: Memoization en Dynamic Programming.
\end{itemize}
}

\resumen{
\begin{itemize}
    \item Space Complexity = cuánta RAM EXTRA usa el algoritmo.
    \item No cuenta el espacio del input original (Auxiliary Space).
    \item Estructuras de datos (arrays, maps) creadas \$\rightarrow\$ O(tamaño).
    \item Recursión \$\rightarrow\$ cada nivel cuenta como O(1) en el Call Stack.
    \item O(1) espacio = "In-place", es lo más eficiente en memoria.
\end{itemize}
}

\end{document}
